\section*{Resumo}
\addcontentsline{toc}{section}{Resumo}

A primeira parte deste projeto propôs uma arquitetura de segurança em camadas para
comunicações MQTT e um plano para a avaliar. Esta segunda parte apresenta o
protótipo que a implementa, as medições feitas num Raspberry~Pi~5 e a validação do
protótipo perante os ataques previstos. Cada camada pode ser ligada e desligada de
forma independente, o que permitiu medir o custo e o benefício de cada uma em doze
configurações (K0 a K11).

O transporte (mTLS) e o envelope AES-GCM custam cada um 7 a 8~$\mu$s de
processador por mensagem no \emph{gateway}, e as restantes verificações, juntas,
cerca de 2~$\mu$s; as assinaturas Ed25519 custam mais 151~$\mu$s e limitam o
débito a 5\,000~mensagens/s. O transporte e o envelope impedem, cada um, quase
todos os ataques de um estranho, mas os de um dispositivo comprometido só são
impedidos pela combinação de todas as camadas, que K8, K9 e K11 conseguem. A
configuração recomendada é K8 para a telemetria (30\,000~mensagens/s, p99 de
0,26~ms) e K9 para as ordens a atuadores, como a primeira parte previa. Das cinco
hipóteses, três confirmam-se; refuta-se que o custo do TLS esteja no
estabelecimento da ligação, e não se confirma que a limitação de tráfego poupe
processador quando a criptografia é AES em \emph{hardware}.

\medskip
\noindent\textbf{Palavras-chave:} MQTT, Internet das Coisas, segurança em camadas,
proteção fim-a-fim, avaliação de desempenho, Raspberry Pi.
